\section{Estado del arte}
\label{sec:mt-estado-arte}

Esta sección revisa los trabajos previos sobre los que se apoya la
propuesta, organizados por familias de enfoques. Se examinan primero
los métodos clásicos de selección de características y, a continuación,
las líneas de optimización cuántica de problemas combinatorios y el uso
de átomos de Rydberg en tareas de optimización. La última subsección se
centra en las formulaciones cuánticas de la selección de
características, que constituyen el antecedente directo del método
estudiado.

\subsection{Métodos clásicos de selección de características}

Sobre las tres familias descritas en el marco teórico
(sección~\ref{sec:mt-fs}), la literatura ha consolidado un repertorio de
selectores que reaparecen de forma
recurrente como referencia. Entre los \emph{filter}, la puntuación F de ANOVA
y la información mutua son los criterios supervisados más utilizados por
su bajo coste, y la varianza se mantiene como criterio no supervisado de
control~\cite{solorio2020}. Entre los métodos \emph{embedded}, los más
extendidos son la regularización $\ell_1$ y las importancias por
impureza de los bosques aleatorios~\cite{breiman2001}, mientras que el
criterio mRMR~\cite{peng2005} sigue siendo la referencia para incorporar
la redundancia de forma explícita. Entre los \emph{wrapper}, el algoritmo
Boruta~\cite{kursa2010}, que compara la importancia de cada variable con
la de copias permutadas (variables \textit{sombra}) y retiene solo las
que la superan de forma consistente, se ha convertido en uno de los
puntos de comparación habituales en la literatura reciente de selección
cuántica, pese a que su coste crece con los reentrenos del modelo base.

Esa misma literatura ha señalado que el rendimiento predictivo no basta
para caracterizar un selector. La estabilidad de los subconjuntos ante
perturbaciones como el cambio de semilla o de muestra se cuantifica con
índices de solape corregidos por azar, como el de
Kuncheva~\cite{kuncheva2007}: un selector inestable produce subconjuntos
poco reproducibles aunque su rendimiento medio sea alto. La estabilidad
y la redundancia interna se han establecido así como dimensiones de
evaluación complementarias al rendimiento.

\subsection{Optimización cuántica de problemas combinatorios}

La resolución de problemas QUBO es una de las aplicaciones más
estudiadas de la computación cuántica actual~\cite{schuld2021,
wittek2014}. Se han explorado tres líneas principales: el
\textit{quantum annealing} o recocido cuántico, un hardware dedicado que
relaja lentamente el sistema hacia el estado fundamental de un
hamiltoniano para hallar el mínimo de la función de coste; los algoritmos
variacionales sobre procesadores de puertas; y la simulación
hamiltoniana analógica sobre plataformas como los átomos neutros. Las
tres comparten el mismo esquema conceptual, codificar la función de
coste en un hamiltoniano y buscar su estado fundamental, y se
diferencian en el mecanismo de búsqueda y en las restricciones que el
hardware impone al problema. Para problemas con restricciones de
exclusión, la vía analógica resulta especialmente natural porque la
propia física implementa la restricción, sin penalizaciones añadidas
en la función de coste.

\subsection{Átomos de Rydberg en optimización}

La demostración experimental de mayor escala de esta vía es la
resolución del conjunto independiente máximo sobre grafos de discos
unitarios con matrices de hasta 289 átomos~\cite{ebadi2022}, que
mostró tanto el potencial de la plataforma como la dependencia del
rendimiento respecto de la dureza de la instancia. En paralelo, una
línea activa de investigación, a la que pertenece el grupo en el que
se enmarca este TFG, estudia cómo diseñar los perfiles de control
para preservar la adiabaticidad: desde el control asistido por
aprendizaje profundo~\cite{ding2021} hasta la optimización de
\textit{drivings} y la caracterización espectral de la dificultad de
instancias MWIS en simuladores de átomos neutros~\cite{arribas2025}.
Estos trabajos delimitan qué tamaños y geometrías de problema son
tratables hoy en simulación y motivan la envolvente operativa que se
define en el capítulo~\ref{cap:metodos}.

\subsection{Selección de características mediante formulaciones cuánticas}

La formulación de la selección de características como problema QUBO,
combinando relevancia e información mutua par a par, fue introducida
por Mücke et al.~\cite{mucke2023}, que la evaluaron sobre hardware
cuántico de puertas y \textit{annealers} con resultados competitivos
frente a heurísticas clásicas en problemas de tamaño moderado. Sobre
átomos neutros, la propuesta de referencia es el método QFS del grupo
en el que se desarrolla este trabajo~\cite{orquin2026}, que adapta esa
misma formulación QUBO al hardware de Rydberg: codifica la relevancia
por información mutua en los campos de control locales y la redundancia
en la disposición espacial de los átomos. En sus experimentos sobre
problemas binarios de hasta una veintena de variables, QFS iguala el
área bajo la curva de Boruta empleando subconjuntos más compactos y, en
uno de los conjuntos, la supera. Un hallazgo relevante de esa
evaluación es que el rendimiento se estabiliza más allá de unas siete
variables, lo que sugiere que Boruta tiende a sobreestimar el número de
características relevantes; el método cuántico ofrece subconjuntos más
parcos sin pérdida apreciable. La propia línea de trabajo contrasta
además las soluciones del simulador con el óptimo de la formulación
QUBO obtenido mediante optimización clásica exacta~\cite{orquin2026},
una práctica de control que la evaluación diseñada en este trabajo
adopta.

El hueco que este TFG aborda se sitúa precisamente ahí: la evidencia
preliminar disponible del método se limita a problemas binarios de
dimensionalidad moderada, evaluados con un único protocolo de
clasificación. Queda pendiente una evaluación contra una referencia
clásica contrastada que cubra también escenarios multiclase,
desbalanceados y de alta dimensionalidad.

La laguna no es únicamente cuantitativa. Para valorar un método como
QFS no basta con registrar si iguala o no al baseline: su función de
coste mezcla relevancia y redundancia, y su implementación física añade
un segundo nivel de dificultad al embebido geométrico y a la dinámica
analógica. Por ello, una evaluación útil debe permitir separar tres
preguntas: qué régimen del dato se está resolviendo, qué familia de
selectores clásicos constituye el espejo apropiado y si un fallo, en
caso de aparecer, procede del criterio o del optimizador. Construir esa
referencia y convertirla en un diagnóstico variable a variable es la
contribución de este trabajo.
